\documentclass[11pt]{article}
\usepackage[margin=1in]{geometry}
\usepackage{amsmath,amssymb,amsthm,booktabs,graphicx,microtype}
\usepackage[hidelinks]{hyperref}

\newtheorem{theorem}{Theorem}
\newtheorem{proposition}[theorem]{Proposition}
\newtheorem{lemma}[theorem]{Lemma}
\newtheorem{corollary}[theorem]{Corollary}
\theoremstyle{remark}
\newtheorem{remark}[theorem]{Remark}
\newcommand{\R}{\mathbb R}
\newcommand{\N}{\mathbb N}
\newcommand{\cL}{\mathcal L}

\title{Critical Vicsek super-Riesz endpoint:\\
one-cut stability and a finite-horizon estimate}
\author{Candidate partial result for arXiv:2606.05475}
\date{12 August 2026}

\begin{document}
\maketitle

\begin{center}
\fbox{\parbox{0.91\textwidth}{\textbf{Status.}
Candidate partial result, likely valid.  This packet does not claim to decide
the full critical inequality.  It proves that all one-edge cut tests are
uniformly harmless at the critical exponent and proves a logarithmic
finite-horizon bound.}}
\end{center}

\section{Source question}

The source is Joseph Feneuil, \emph{Reverse inequalities for super-Riesz
transforms on graphs with a slow diffusion}, arXiv:2606.05475v2
\cite{Feneuil}.  On a $D$-dimensional Vicsek graph it proves that
\[
 (R_{p,\gamma})\qquad
 \|\,|\nabla f|\,\|_p\leq C\|\Delta^\gamma f\|_p
\]
holds below
\[
 \gamma^*(p)=\frac1{D+1}+\frac{D-1}{(D+1)p}
\]
and fails above it, leaving equality open.  The reverse inequality is
\[
 (RR_{p,\gamma})\qquad
 \|\Delta^\gamma f\|_p\leq C\|\,|\nabla f|\,\|_p.
\]
On a tree, the source's primitive and duality arguments pair the critical
Riesz inequality at $p$ with the critical reverse inequality at $p'$.  Thus
understanding either family at equality is directly relevant to the question.

Figure~\ref{fig:question} reproduces the abstract's exact statement that the
critical exponent remains open.

\begin{figure}[ht]
\centering
\includegraphics[width=\textwidth]{figures/open_problem_crop.png}
\caption{The open endpoint statement on page 1 of arXiv:2606.05475v2.}
\label{fig:question}
\end{figure}

\section{Setting and notation}

Let $(G,\mu)$ be an infinite weighted tree satisfying the assumptions used in
the source: uniform laziness/non-degeneracy (LB), $D$-Ahlfors regularity, and
the sub-Gaussian upper estimate $(UE)_\beta$.  Thus
\[
 V(x,r)\asymp (1+r)^D
\]
and, with harmless indexing changes,
\[
 p_k(x,y)\leq \frac{C}{V(x,k^{1/\beta})}
 \exp\!\left[-c\left(\frac{d(x,y)^\beta}{k}\right)^{1/(\beta-1)}\right].
 \tag{2.1}
\]
For a Vicsek graph $\beta=D+1$.  We write $P$ for the Markov operator and
$\Delta=I-P$.

If $e=\{a,b\}$ is a non-loop edge, deleting $e$ gives two components.  Let
$A_e$ be either one and put $h_e=\mathbf 1_{A_e}$.  Finally set
\[
 \rho_e(x)=1+\min\{d(x,a),d(x,b)\}.
\]

\section{Idea of the proof}

A one-edge cut has the smallest possible gradient: it is supported only at
the cut.  Fractional powers are nonlocal, however, so this test could in
principle still destroy the endpoint.  The binomial formula for
$(I-P)^\gamma$ turns its value at $x$ into a weighted sum of probabilities
that the random walk crosses the cut.  Before time comparable with
$\rho_e(x)^\beta$, the sub-Gaussian estimate makes that probability
exponentially small; afterward the fractional coefficient tail is exactly
$\rho_e(x)^{-\beta\gamma}$.  Ahlfors shell summation then gives the sharp
integrability threshold $p\beta\gamma>D$.

At the critical exponent, the source's Poincar\'e exponent is exactly
$s=\beta\gamma^*(p)$.  Its tent-space lemma remains valid with the local
exponent $q=p$; what fails in the published strong-type proof is summation
over infinitely many diffusion times.  Truncating at $K$ and applying
Minkowski in $L^{p/2}$ costs only the square root of the harmonic sum
$\sum_{k\leq K}k^{-1}$.

The full endpoint remains open because neither argument controls coherent
superpositions of cuts over unboundedly many self-similar scales.

\section{Uniform stability of one-edge cuts}

\begin{theorem}[one-edge cut estimate]\label{thm:cut}
Let $0<\gamma<1$.  Uniformly in the edge $e$ and in the choice of component,
\[
 |\Delta^\gamma h_e(x)|\leq C_\gamma\rho_e(x)^{-\beta\gamma}
 \qquad (x\in G).
 \tag{4.1}
\]
Consequently
\[
 \Delta^\gamma h_e\in L^p(G)
 \quad\text{and}\quad
 \|\Delta^\gamma h_e\|_p\leq C_{p,\gamma}
 \tag{4.2}
\]
whenever $p\beta\gamma>D$.  Moreover
$\|\,|\nabla h_e|\,\|_p\asymp1$, uniformly in $e$.
\end{theorem}

\begin{proof}
For $0<\gamma<1$, the absolutely convergent scalar binomial identity on
$[-1,1]$ is
\[
 (1-z)^\gamma=\sum_{k=1}^\infty b_k(1-z^k),\qquad
 b_k=(-1)^{k+1}\binom{\gamma}{k}>0,
 \quad b_k\asymp_\gamma k^{-1-\gamma}.
 \tag{4.3}
\]
It applies pointwise to the bounded function $h_e$ and gives
\[
 \Delta^\gamma h_e(x)=\sum_{k=1}^\infty
 b_k\bigl(h_e(x)-P^kh_e(x)\bigr).
 \tag{4.4}
\]

Suppose first that $x\in A_e$.  Since every path from $x$ to $A_e^c$ crosses
$e$, every $y\in A_e^c$ has $d(x,y)\geq \rho_e(x)-1$.  Annular summation in
(2.1), using Ahlfors regularity, gives constants $C,c'>0$ such that
\[
 0\leq 1-P^kh_e(x)
 =\sum_{y\in A_e^c}p_k(x,y)m(y)
 \leq C\exp\!\left[-c'
 \left(\frac{\rho_e(x)^\beta}{k}\right)^{1/(\beta-1)}\right]
 \tag{4.5}
\]
when $k\leq\rho_e(x)^\beta$, while the trivial upper bound is $1$ for all
$k$.  The same argument, with the two components exchanged, applies when
$x\notin A_e$.

Write $R=\rho_e(x)^\beta$.  Equations (4.3)--(4.5), split at $k=R$, imply
\begin{align*}
 |\Delta^\gamma h_e(x)|
 &\leq C\sum_{k\leq R} k^{-1-\gamma}
     e^{-c'(R/k)^{1/(\beta-1)}}
   +C\sum_{k>R}k^{-1-\gamma}\\
 &\leq C R^{-\gamma}=C\rho_e(x)^{-\beta\gamma}.
\end{align*}
The first sum is estimated by the change of scale $u=k/R$ (or by dyadic
blocks), and the second is an integral tail.  This proves (4.1).

For $j\geq0$, Ahlfors regularity bounds the measure of
$\{x:2^j\leq\rho_e(x)<2^{j+1}\}$ by $C2^{jD}$.  Therefore
\[
 \|\Delta^\gamma h_e\|_p^p
 \leq C\sum_{j\geq0}2^{j(D-p\beta\gamma)},
\]
which converges precisely under the asserted strict inequality.  Finally,
$\nabla h_e$ vanishes away from the two orientations of $e$; (LB), Ahlfors
regularity at bounded radius, and uniform local finiteness make its $L^p$
norm comparable with $1$.
\end{proof}

\begin{corollary}[critical cut stability]\label{cor:critical}
On a $D$-dimensional Vicsek graph, for every $1<p<\infty$,
\[
 \sup_e\|\Delta^{\gamma^*(p)}h_e\|_p<\infty.
\]
Indeed,
\[
 p\beta\gamma^*(p)=p+D-1>D.
\]
Thus a counterexample at equality cannot consist of a single bottleneck cut.
\end{corollary}

\section{A logarithmic finite-horizon endpoint bound}

For $K\geq1$, define the truncated Lusin area functional
\[
 \cL_{\gamma,K}f(x)^2=
 \sum_{k=1}^K\ \sum_{d(x,y)^\beta<k}
 \frac{k^{1-2\gamma}m(y)}{V(x,k^{1/\beta})}
 |\Delta P^{k-1}f(y)|^2.
 \tag{5.1}
\]
The source proves that the corresponding full functional is equivalent in
$L^p$ to $\Delta^\gamma f$ (via its horizontal and conical Littlewood--Paley
theorems).

\begin{theorem}[finite-horizon critical estimate]\label{thm:log}
Let $2\leq p<\infty$, $\beta=D+1$, and
$\gamma=\gamma^*(p)$.  For every compactly supported $f$ and every $K\geq2$,
\[
 \|\cL_{\gamma,K}f\|_p
 \leq C_p\sqrt{1+\log K}\,\|\,|\nabla f|\,\|_p.
 \tag{5.2}
\]
The constant is independent of $K$ and $f$.
\end{theorem}

\begin{proof}
Put $s=1+(D-1)/p=\beta\gamma$.  The source's Vicsek Poincar\'e theorem gives
$(P_p)_s$.  Its area-functional estimate preceding the proof of the main
reverse theorem is valid with local exponent $q=p$; the later restriction
$q<p$ is used only to obtain a strong maximal-function estimate.  Apply that
area estimate to the collection $f_k=f$ for $k\leq K$ and $f_k=0$ afterward.
Because $s/\beta=\gamma$, it gives
\[
 \|\cL_{\gamma,K}f\|_p
 \leq C_p\left\|\mathcal A^p
  (\{\mathbf1_{k\leq K}|\nabla f|\}_{k\geq1})\right\|_p.
 \tag{5.3}
\]

Set $g=|\nabla f|$ and $r_k=k^{1/\beta}$.  From the definition of
$\mathcal A^p$, its square has the form
\[
 \sum_{k=1}^K\frac1k H_k(x),
\quad
 H_k(x)=\frac1{V(x,r_k)}\sum_{y\in B(x,r_k)}m(y)
 \left(\frac1{V(y,r_k)}\sum_{z\in B(y,r_k)}|g(z)|^p m(z)
 \right)^{2/p}.
 \tag{5.4}
\]
Doubling, Jensen, and Fubini imply
\[
 \|H_k\|_{p/2}\leq C\|g\|_p^2
 \tag{5.5}
\]
uniformly in $k$.  (For $p=2$ this is just preservation of mass up to the
doubling constant; for $p>2$, apply Jensen to the exponent $2/p\leq1$ and
then use the same Fubini estimate.)  Since $p/2\geq1$, Minkowski yields
\begin{align*}
 \|\mathcal A^p(\{\mathbf1_{k\leq K}g\})\|_p^2
 &=\left\|\sum_{k=1}^K k^{-1}H_k\right\|_{p/2}\\
 &\leq\sum_{k=1}^K k^{-1}\|H_k\|_{p/2}
 \leq C(1+\log K)\|g\|_p^2.
\end{align*}
Combining this with (5.3) proves (5.2).
\end{proof}

\begin{remark}
At $p=2$, the exact identity
$\|\Delta^{1/2}f\|_2=\|\,|\nabla f|\,\|_2$ removes the logarithm.  For
$p>2$, Theorem~\ref{thm:log} is not claimed sharp.  Since
$\cL_{\gamma,K}\uparrow\cL_{\gamma,\infty}$, a uniform version of (5.2)
would prove the reverse critical inequality.  The theorem shows precisely
that the published method loses no power of the spatial radius: its remaining
loss is harmonic accumulation over diffusion scales.
\end{remark}

\section{Computational stress test and full-upgrade attempts}

The canonical finite Vicsek tree was built recursively from five copies of an
X-tree.  With the lazy normalized Laplacian, the induced norm of
$\Delta^{\gamma^*(4)}T:\ell^4(E)\to L^4(V)$ was optimized, where $T$ is the
rooted tree primitive.  Dense diagonalization was used through level four;
levels five and six used the Balakrishnan integral with sparse shifted solves.

\begin{center}
\begin{tabular}{@{}cccc@{}}
\toprule
level $n$ & vertices & critical reverse norm & norm$/ (n+1)^{1/4}$\\
\midrule
0 & 5     & 0.8907 & 0.891\\
1 & 21    & 1.0861 & 0.913\\
2 & 101   & 1.2299 & 0.935\\
3 & 501   & 1.3361 & 0.945\\
4 & 2501  & 1.41905& 0.949\\
5 & 12501 & 1.49025& 0.952\\
6 & 62501 & 1.54647& 0.951\\
\bottomrule
\end{tabular}
\end{center}

The normalization is strikingly stable and predicts the multiscale law
$(n+1)^{1/2-1/p}$ at $p=4$.  The conjugate and opposite-direction probes show
the analogous growth.  A fixed level-four witness embedded into levels five
and six stabilizes instead; growth appears only after re-optimization.  Hence
the effect is not a boundary tail of one local cut but a hierarchy that changes
with the number of available scales.

This is not a proof of divergence.  Eight focused upgrade attempts were made:
(1) audit the endpoint in the source proof; (2) dense finite approximants;
(3) Hardy/BMO interpolation; (4) comparison above and below the line;
(5) critical Besov versus square-function analysis; (6) hierarchical
extremizer decomposition; (7) sparse tail embeddings; and (8) sparse
optimization through level six.  The remaining credible route is a uniform
cell-block lemma giving matching infinite-volume lower bounds.  It was not
proved here.

\section{Verification, limitations, and review recommendation}

\begin{itemize}
\item Theorem~\ref{thm:cut} uses only (LB), Ahlfors regularity, $(UE)_\beta$,
and the tree cut property, so it applies to the source's broadened definition
of a Vicsek graph, not only the self-similar example.
\item The strict condition $p\beta\gamma>D$ is essential to the shell argument.
At the source endpoint it has margin $p-1$, so no equality ambiguity is hidden
in Corollary~\ref{cor:critical}.
\item Theorem~\ref{thm:log} is proved only for $p\geq2$, where Minkowski in
$L^{p/2}$ is available.  It is a truncated area-functional statement, not the
full reverse inequality.
\item The computations in \texttt{code/} are reproducibility evidence only.
They use finite Neumann-type approximants and floating-point quadrature; no
numerical value is used in either proof.
\item Cheap run indexes and the locally ingested corpus were searched for the
arXiv id, the critical formula, and Vicsek endpoint terminology.  No duplicate
result packet or later full answer was found.  This was not an exhaustive
external novelty search, so novelty confidence is moderate.
\end{itemize}

\noindent\textbf{Human-review recommendation.}
Promote the two proved estimates as a substantial partial result after checking
the use of the $q=p$ instance of the source's area lemma and the Poisson-kernel
annular summation in (4.5).  Do not promote the numerical full-failure
conjecture without the missing uniform cell-block lower bound.

\begin{thebibliography}{9}
\bibitem{Feneuil}
J. Feneuil,
\emph{Reverse inequalities for super-Riesz transforms on graphs with a slow
diffusion}, arXiv:2606.05475v2, 2026.
\end{thebibliography}

\end{document}
